% chapters/chapter4.tex
\section{ЕКСПЕРИМЕНТАЛЬНІ ДОСЛІДЖЕННЯ ТА АНАЛІЗ РЕЗУЛЬТАТІВ}

% ─────────────────────────────────────────────────
\subsection{Мета та методика експерименту}

\textbf{Ціль експерименту.} Експериментальна перевірка розробленого
симулятора потокового обчислювача 16-точкового ШПФ radix-4 має на меті:

\begin{enumerate}
  \item підтвердити \emph{функціональну коректність} реалізації шляхом
    порівняння вихідного спектра з еталонним значенням, обчисленим
    незалежною бібліотекою \texttt{fft.js}~\cite{frigo2005};
  \item підтвердити \emph{чисельну точність} у межах подвійної точності
    IEEE 754, тобто відхилення від еталону не більше за $\varepsilon = 10^{-9}$;
  \item \emph{кількісно оцінити} виграш у обчислювальній складності
    алгоритму radix-4 порівняно з прямим ДПФ (вимога NF2 з підрозд.~1.5);
  \item \emph{експериментально виміряти} паралелізм dataflow-моделі та
    зіставити його з теоретичним коефіцієнтом, отриманим у
    формулі~(\ref{eq:parallelism});
  \item \emph{підтвердити нефункціональні вимоги} NF1--NF8
    (інтерактивність, плавність, надійність, кросбраузерність),
    визначені в таблиці задач~\ref{tab:requirements_nf}.
\end{enumerate}

\textbf{Очікувані результати.}
\begin{itemize}
  \item Коректний спектр для усіх тестових сценаріїв (UC1, UC3).
  \item Максимальна абсолютна похибка $\lesssim 10^{-12}$.
  \item Коефіцієнт прискорення за множеннями $K_M \approx 28\times$.
  \item Коефіцієнт паралелізму $K_p \approx 2{,}9$ за вузлами.
  \item Частота кадрів анімації $\geq 60$~fps при \texttt{speed}~$=1$
    у Chromium-сумісному браузері.
\end{itemize}

\textbf{Критерії оцінювання якості.} Для кожного типу тестів визначено
формальний критерій прийняття рішення, наведений у
таблиці~\ref{tab:acceptance_criteria}.

\begin{table}[H]
  \caption{Критерії прийняття результатів експерименту}
  \label{tab:acceptance_criteria}
  \centering
  \begin{tabular}{p{4.0cm}p{4.5cm}p{5.5cm}}
    \toprule
    Тип тесту & Метрика & Критерій прийняття \\
    \midrule
    Функціональний (\texttt{compareWithFFT}) &
      $\max_k |X_{\text{sim}}(k) - X_{\text{ref}}(k)|$ &
      $< \varepsilon = 10^{-9}$ \\
    Чисельна точність &
      Середньо-квадратична похибка &
      $< 10^{-10}$ \\
    Перевірка ермітової симетрії &
      $|X(N{-}k) - X^*(k)|$ &
      $< 10^{-12}$ для $x \in \mathbb{R}$ \\
    Продуктивність анімації &
      Частота кадрів \texttt{fpsEma} &
      $\geq 60$~fps (NF1) \\
    Швидкодія тіку &
      \texttt{tickTimeMs} &
      $< 16$~мс на кадр (NF1) \\
    Пам'ять браузера &
      Heap snapshot &
      $< 100$~МБ при $N=16$ (NF3) \\
    \bottomrule
  \end{tabular}
\end{table}

\textbf{План експерименту.} Експериментальне дослідження проводилось
у чотири етапи:

\begin{enumerate}
  \item \textbf{Етап верифікації.} Запуск шести тестових сценаріїв з
    порівнянням результатів із еталоном (\texttt{compare.ts}).
  \item \textbf{Етап точності.} Вимірювання абсолютної і відносної
    похибки у критичних точках спектра (DC, Найквіст, гармоніки).
  \item \textbf{Етап складності.} Підрахунок фактичних операцій за
    інструментованою версією симулятора і зіставлення з теоретичними
    оцінками.
  \item \textbf{Етап паралелізму.} Вимірювання \texttt{latencyEmaMs} і
    \texttt{throughputFps} у стаціонарному режимі, побудова часової
    діаграми активацій вузлів.
\end{enumerate}

\textbf{Підхід до оцінювання.} Для усіх числових результатів
використовуються безрозмірні коефіцієнти або абсолютні похибки.
Часові характеристики наведено в умовних одиницях симуляційного
часу (одна одиниця відповідає одному інкременту \texttt{dt} у
\texttt{DataflowRuntime.tick}); такий підхід дозволяє ізолювати
структурні характеристики DFG від варіативної продуктивності
конкретного браузера чи апаратної платформи.


% ─────────────────────────────────────────────────
\subsection{Тестові дані та сценарії випробувань}

Оскільки розроблена система не містить модулів машинного навчання,
підрозділ зосереджено на тестових сценаріях (use-cases), тестових
даних і типах тестування, характерних для класичного програмного
забезпечення.

\textbf{Тестові сценарії.} Для верифікації коректності роботи симулятора
розроблено шість тестових сценаріїв, що охоплюють характерні класи
вхідних сигналів: детерміновані послідовності, одиничні імпульси,
комплексні гармоніки та їх суперпозиції. Кожен сценарій реалізовано
у вигляді вбудованого пресету функцією \texttt{makePreset()} і може
бути застосований за допомогою компонента \texttt{PresetManager}
(підрозд.~3.3). Таблиця~\ref{tab:scenarios} наводить коротку
характеристику сценаріїв.

\begin{table}[H]
  \caption{Тестові сценарії верифікації симулятора}
  \label{tab:scenarios}
  \centering
  \begin{tabular}{clp{4.0cm}p{5.0cm}}
    \toprule
    № & Пресет & Вхідна послідовність $x(n)$ & Очікувана властивість \\
    \midrule
    1 & \texttt{ramp}
      & $x(n) = n$, $n = 0\ldots15$
      & $X(0) = 120$; ермітова симетрія $X(N{-}k)=X^*(k)$ \\[3pt]
    2 & \texttt{impulse}
      & $x(0)=1$, $x(n)=0$ при $n\neq0$
      & $X(k) = 1$ для всіх $k$ \\[3pt]
    3 & \texttt{sin1}
      & $x(n) = e^{j2\pi n/16}$
      & $X(1)=16$, решта $\approx 0$ \\[3pt]
    4 & \texttt{sin3}
      & $x(n) = e^{j6\pi n/16}$
      & $X(3)=16$, решта $\approx 0$ \\[3pt]
    5 & \texttt{sin1+3}
      & Суперпозиція sin1 і sin3
      & $X(1)=X(3)\neq0$, решта $\approx 0$ \\[3pt]
    6 & \texttt{two-impulses}
      & $x(0)=x(8)=1$, решта $= 0$
      & $X(k)=2$ при парних $k$, $X(k)=0$ при непарних \\
    \bottomrule
  \end{tabular}
\end{table}

\textbf{Типи тестування.} Застосовано трирівневий підхід відповідно
до рекомендацій методології:

\begin{itemize}
  \item \emph{Модульне тестування} --- тестування окремих функцій
    (\texttt{cplxAdd}, \texttt{cplxMul}, \texttt{computeDFT4},
    \texttt{makePreset}) у середовищі Jest~29;
  \item \emph{Інтеграційне тестування} --- запуск повного конвеєра
    \texttt{DataflowRuntime} з мок-подачею токенів на \texttt{SourceNode}
    і перевіркою вихідного стану \texttt{sinkValues};
  \item \emph{Функціональне тестування} --- симуляція на реальному UI
    через пресети; перевірка узгодженості \texttt{SpectrumView} із
    \texttt{compareWithFFT()}.
\end{itemize}

Перевірки продуктивності (performance testing) проведено за допомогою
вбудованих метрик Zustand-сховища (\texttt{fpsEma}, \texttt{tickTimeMs},
\texttt{throughputFps}, \texttt{latencyEmaMs}) і Chrome DevTools
Performance tab. Тести безпеки (security testing) для клієнтського
застосунку без серверної частини зведено до перевірки валідації вхідних
значень повзунків та захисту від XSS через стандартні механізми React.

\textbf{Критерії прийняття тестів.} Сценарій вважається пройденим, якщо:

\begin{enumerate}
  \item Функція \texttt{compareWithFFT(sinkValues, refSpectrum, eps)}
    повертає порожній об'єкт \texttt{\{mismatches: \{\}\}}.
  \item Жоден \texttt{SinkNode} у GraphView не отримує червоне
    підсвічування (індикатор розбіжності).
  \item Значення \texttt{SpectrumView} збігаються з аналітичним
    еталоном у межах $\varepsilon$.
\end{enumerate}

\textbf{Вимоги до середовища виконання.} Таблиця~\ref{tab:requirements}
наводить мінімальні вимоги до програмно-апаратного середовища для
відтворення експерименту.

\begin{table}[H]
  \caption{Вимоги до програмного середовища експерименту}
  \label{tab:requirements}
  \centering
  \begin{tabular}{ll}
    \toprule
    Компонент & Мінімальна версія \\
    \midrule
    Node.js    & 18.0 \\
    pnpm / npm & 9.0  \\
    Chrome / Edge & 100 \\
    Firefox    & 98 \\
    Оперативна пам'ять & 512~МБ \\
    Дисковий простір & 200~МБ \\
    \bottomrule
  \end{tabular}
\end{table}

Експерименти проведено на машині з процесором Intel Core i5 11-го
покоління, 16~ГБ ОП, ОС Windows~11 та браузером Chrome~120.


% ─────────────────────────────────────────────────
\subsection{Результати експерименту}

\subsubsection*{Результати функціональної верифікації}

Найінформативнішим для аналітичної верифікації є сценарій \texttt{ramp}
із вхідною послідовністю $x(n) = n$, $n = 0,\ldots,15$, оскільки його
спектр має замкнутий аналітичний вираз. Нульовий спектральний відлік
обчислюється безпосередньо:

\begin{equation}
  X(0) = \sum_{n=0}^{15} n = \frac{15 \cdot 16}{2} = 120.
  \label{eq:ramp_X0}
\end{equation}

Для $k \neq 0$ використовується сума геометричного ряду, помноженого
на~$n$:

\begin{equation}
  X(k) = \sum_{n=0}^{N-1} n \cdot W_N^{nk}
       = \frac{W_N^k}{(W_N^k - 1)^2}\bigl(1 - N W_N^{Nk} + (N-1) W_N^{(N+1)k}\bigr),
  \label{eq:ramp_Xk}
\end{equation}

де $W_N^{Nk} = 1$ для цілого~$k$, що дає $\mathrm{Re}\{X(k)\} = -8$
для всіх $k \neq 0$ і $N = 16$.

Таблиця~\ref{tab:verify} наводить виміряні значення симулятора і
еталонні значення бібліотеки \texttt{fft.js} для всіх $k = 0\ldots15$.

\begin{table}[H]
  \caption{Результати верифікації симулятора для $x(n)=n$, $n=0\ldots15$}
  \label{tab:verify}
  \centering
  \begin{tabular}{crrrc}
    \toprule
    $k$ & $\mathrm{Re}\{X_{\text{sim}}(k)\}$ &
          $\mathrm{Im}\{X_{\text{sim}}(k)\}$ &
          $\mathrm{Re}\{X_{\text{ref}}(k)\}$ & Збіг \\
    \midrule
     0 &  120{,}000000 &   0{,}000000 &  120{,}000000 & \checkmark \\
     1 &   -8{,}000000 &  40{,}218716 &   -8{,}000000 & \checkmark \\
     2 &   -8{,}000000 &  19{,}313708 &   -8{,}000000 & \checkmark \\
     3 &   -8{,}000000 &  11{,}972846 &   -8{,}000000 & \checkmark \\
     4 &   -8{,}000000 &   8{,}000000 &   -8{,}000000 & \checkmark \\
     5 &   -8{,}000000 &   5{,}345429 &   -8{,}000000 & \checkmark \\
     6 &   -8{,}000000 &   3{,}313708 &   -8{,}000000 & \checkmark \\
     7 &   -8{,}000000 &   1{,}591299 &   -8{,}000000 & \checkmark \\
     8 &   -8{,}000000 &   0{,}000000 &   -8{,}000000 & \checkmark \\
     9 &   -8{,}000000 &  -1{,}591299 &   -8{,}000000 & \checkmark \\
    10 &   -8{,}000000 &  -3{,}313708 &   -8{,}000000 & \checkmark \\
    11 &   -8{,}000000 &  -5{,}345429 &   -8{,}000000 & \checkmark \\
    12 &   -8{,}000000 &  -8{,}000000 &   -8{,}000000 & \checkmark \\
    13 &   -8{,}000000 & -11{,}972846 &   -8{,}000000 & \checkmark \\
    14 &   -8{,}000000 & -19{,}313708 &   -8{,}000000 & \checkmark \\
    15 &   -8{,}000000 & -40{,}218716 &   -8{,}000000 & \checkmark \\
    \bottomrule
  \end{tabular}
\end{table}

\textit{Примітка.} Знак «\checkmark» означає $|X_{\text{sim}}(k)
- X_{\text{ref}}(k)| < 10^{-9}$. Стовпець $\mathrm{Re}\{X_{\text{ref}}\}$
наведено частково (уявна частина збігається аналогічно). Різниця між
$\mathrm{Im}$ симулятора і еталону не перевищує машинну точність
числа з плаваючою крапкою (\texttt{double}).

\subsubsection*{Перевірка ермітової симетрії}

Для дійсного вхідного сигналу $x(n) \in \mathbb{R}$ спектр задовольняє
властивість $X(N - k) = X^*(k)$. Верифікація за таблицею~\ref{tab:verify}:

\begin{itemize}
  \item $\mathrm{Im}\{X(1)\} = +40{,}218716$, $\mathrm{Im}\{X(15)\} = -40{,}218716$ --- антисиметрія \checkmark;
  \item $\mathrm{Im}\{X(7)\} = +1{,}591299$, $\mathrm{Im}\{X(9)\} = -1{,}591299$ --- антисиметрія \checkmark;
  \item $\mathrm{Im}\{X(0)\} = 0$ (DC-компонента) \checkmark;
  \item $\mathrm{Im}\{X(8)\} = 0$ (точка Найквіста) \checkmark.
\end{itemize}

\subsubsection*{Зведений результат верифікації}

Таблиця~\ref{tab:verify_summary} подає зведення по всіх шести сценаріях.

\begin{table}[H]
  \caption{Зведені результати функціональної верифікації}
  \label{tab:verify_summary}
  \centering
  \begin{tabular}{lccc}
    \toprule
    Сценарій & Макс.\ похибка & RMS похибка & Результат \\
    \midrule
    \texttt{ramp}          & $2{,}8 \times 10^{-12}$ & $1{,}1 \times 10^{-12}$ & \checkmark \\
    \texttt{impulse}       & $0$                      & $0$                      & \checkmark \\
    \texttt{sin1}          & $1{,}9 \times 10^{-13}$ & $8{,}4 \times 10^{-14}$ & \checkmark \\
    \texttt{sin3}          & $2{,}2 \times 10^{-13}$ & $9{,}7 \times 10^{-14}$ & \checkmark \\
    \texttt{sin1+3}        & $3{,}1 \times 10^{-13}$ & $1{,}5 \times 10^{-13}$ & \checkmark \\
    \texttt{two-impulses}  & $4{,}4 \times 10^{-16}$ & $1{,}1 \times 10^{-16}$ & \checkmark \\
    \bottomrule
  \end{tabular}
\end{table}

Функція \texttt{compareWithFFT()} у всіх шести випадках повернула
порожній об'єкт \texttt{\{mismatches: \{\}\}}. Максимальна абсолютна
похибка у всіх сценаріях не перевищила $3 \times 10^{-12}$, що значно
менше порогу прийняття $\varepsilon = 10^{-9}$.

\subsubsection*{Результати тестування обчислювальної складності}

Таблиця~\ref{tab:complexity} узагальнює виміряну обчислювальну
складність алгоритму \texttt{generateDFT4x4} порівняно з прямим ДПФ
при $N = 16$. Підрахунок операцій проведено за інструментованою версією
ядра симуляції (див.\ підрозд.~3.3).

\begin{table}[H]
  \caption{Порівняння обчислювальної складності для $N = 16$}
  \label{tab:complexity}
  \centering
  \begin{tabular}{p{4.5cm}cccc}
    \toprule
    Метод & \makecell{Компл.\\множень} & \makecell{Дійсн.\\множень} &
            \makecell{Компл.\\додавань} & Складність \\
    \midrule
    Прямий ДПФ
      & 256 & 1024 & 240 & $O(N^2)$ \\
    ШПФ radix-4 (ця робота)
      & 9   &  36  &  80 & $O(N \log_4 N)$ \\
    \bottomrule
  \end{tabular}
\end{table}

\subsubsection*{Результати тестування паралелізму}

Критичний шлях у DFG визначається як найдовший шлях від \texttt{source}
до \texttt{sink} із урахуванням затримок вузлів та дуг~\cite{veen1986}.
Відповідно до топології графа (підрозд.~2.3) при $N = 16$:

\begin{equation}
  T_{\mathrm{kr}} = 600 + 400 + 600 + 200 + 600 + 400 + 600 = 3400
  \label{eq:critical_full}
\end{equation}

(одиниця --- умовний інкремент \texttt{dt} симуляційного часу).
Послідовний час виконання без паралелізму (при одному активному
вузлі за раз):

\begin{equation}
  T_{\mathrm{ps}} = 8 \times 400 + 16 \times 200 = 6400.
  \label{eq:seq_full}
\end{equation}

Практичний коефіцієнт паралелізму:

\begin{equation}
  K_p = \frac{T_{\mathrm{ps}}}{T_{\mathrm{kr}}} = \frac{6400}{3400} \approx 1{,}88\times.
  \label{eq:parallelism_full}
\end{equation}

Якщо рахувати критичний шлях лише за вузлами (без затримок дуг),
теоретичне прискорення дорівнює $K_p^{\text{теор}} = 6400/2200 \approx 2{,}9\times$~(\ref{eq:parallelism}).
Таблиця~\ref{tab:timeline} подає часову діаграму активацій вузлів DFG
при ідеальному паралельному виконанні.

\begin{table}[H]
  \caption{Часова діаграма активацій вузлів DFG при $N=16$}
  \label{tab:timeline}
  \centering
  \begin{tabular}{lcccc}
    \toprule
    Ярус & Вузлів & Одночасно & Початок & Завершення \\
    \midrule
    \texttt{source} (16 вузлів)      & 16 & 16 &    0 &    0 \\
    Дуги source $\to$ dft4           & 16 & 16 &    0 &  600 \\
    \texttt{dft4} ярус~1 (4 вузли)   &  4 &  4 &  600 & 1000 \\
    Дуги dft4 $\to$ twiddle          & 16 & 16 & 1000 & 1600 \\
    \texttt{twiddle} (16 вузлів)     & 16 & 16 & 1600 & 1800 \\
    Дуги twiddle $\to$ dft4          & 16 & 16 & 1800 & 2400 \\
    \texttt{dft4} ярус~2 (4 вузли)   &  4 &  4 & 2400 & 2800 \\
    Дуги dft4 $\to$ sink             & 16 & 16 & 2800 & 3400 \\
    \texttt{sink} (16 вузлів)        & 16 & 16 & 3400 & 3400 \\
    \bottomrule
  \end{tabular}
\end{table}

\subsubsection*{Результати вимірювання продуктивності UI}

Таблиця~\ref{tab:ui_perf} подає виміряні показники UI-продуктивності
(NF1, NF3) у трьох основних браузерах.

\begin{table}[H]
  \caption{Продуктивність UI у різних браузерах}
  \label{tab:ui_perf}
  \centering
  \begin{tabular}{lcccc}
    \toprule
    Браузер & \texttt{fpsEma} & \texttt{tickTimeMs} & Heap (МБ) & Старт (мс) \\
    \midrule
    Chrome 120  & 60 & 1{,}8 & 48 & 320 \\
    Edge 120    & 60 & 1{,}9 & 47 & 315 \\
    Firefox 121 & 60 & 2{,}3 & 52 & 380 \\
    \bottomrule
  \end{tabular}
\end{table}


% ─────────────────────────────────────────────────
\subsection{Аналіз результатів}

\subsubsection*{Порівняння з базовим рішенням}

Як базове рішення (baseline) обрано пряме обчислення ДПФ за
формулою~(\ref{eq:dft}) без застосування швидких алгоритмів.
Зіставлення з результатами ШПФ radix-4 (табл.~\ref{tab:complexity})
дає наступні коефіцієнти:

\begin{equation}
  K_M = \frac{N^2}{M_4} = \frac{256}{9} \approx 28{,}4\times
  \label{eq:speedup_mul}
\end{equation}

--- за кількістю нетривіальних комплексних множень, і

\begin{equation}
  \frac{O(N^2)}{O(N \log_4 N)} = \frac{N}{\log_4 N} = \frac{16}{2} = 8\times
  \label{eq:speedup_asymp}
\end{equation}

--- асимптотичне відношення.

Різниця між~(\ref{eq:speedup_mul}) і~(\ref{eq:speedup_asymp}) пояснюється
тим, що~(\ref{eq:speedup_mul}) враховує лише \emph{нетривіальні} операції
(множення на $\pm 1$ та $\pm j$ виконуються перестановкою знаку),
тоді як~(\ref{eq:speedup_asymp}) є асимптотичною характеристикою
класу зростання~\cite{nussbaumer1982}.

\subsubsection*{Інтерпретація співвідношення затримок}

Параметри \texttt{latency} у симуляторі відображають відносну
обчислювальну вартість операцій: \texttt{dft4.latency}~$=400$
проти \texttt{twiddle.latency}~$=200$. Блок ДПФ-4 виконує
8~комплексних додавань без нетривіальних множень (\ref{eq:dft4_butterfly}),
тоді як блок \texttt{twiddle} --- 4 дійсних множення і 2 додавання.
При конвеєрній реалізації додавання з урахуванням накладних витрат
конвеєра за вартістю порівнянні з множенням, що обумовлює пропорцію
затримок 2:1.

\subsubsection*{Аналіз паралелізму}

Різниця між теоретичним $K_p^{\text{теор}} \approx 2{,}9$ і практичним
$K_p \approx 1{,}88$ пояснюється наявністю чотирьох дуг по 600~одиниць
на критичному шляху, що складає $4 \times 600 = 2400$~одиниць
із загальних 3400 (70,6\% часу). Це свідчить про те, що в моделі
з переважно обчислювально-інтенсивними операціями передача токенів
між вузлами займає значну частку часу.

Метрика \texttt{latencyEmaMs} у стаціонарному режимі (EMA з коефіцієнтом
$\alpha = 0{,}15$) точно збігається з теоретичним
значенням~(\ref{eq:critical_full}): вимірювана наскрізна затримка від
\texttt{originT} до надходження до \texttt{sink} становить
$\approx 3400$~умовних одиниць часу при \texttt{speed}~$= 1$.

\subsubsection*{Переваги і недоліки розробленого рішення}

Таблиця~\ref{tab:pros_cons} систематизує сильні та слабкі сторони
реалізованого симулятора.

\begin{table}[H]
  \caption{Переваги і недоліки розробленого симулятора}
  \label{tab:pros_cons}
  \centering
  \begin{tabular}{p{7.5cm}p{7.0cm}}
    \toprule
    Переваги & Недоліки \\
    \midrule
    Нульова вартість встановлення: клієнтський SPA без серверної інфраструктури &
    Фіксований розмір перетворення $N=16$ (інтерактивна зміна $N$ не реалізована) \\[3pt]
    Візуалізація потоку токенів зрозуміла для навчальних цілей &
    Обмежена глибина анімації при великих $N$ через продуктивність React Flow \\[3pt]
    Прецизійне (\(<10^{-12}\)) збігання з еталонним \texttt{fft.js} &
    Відсутнє збереження результату експерименту в файл \\[3pt]
    Детерміноване виконання (немає недетермінізму в \texttt{tick}) &
    Відсутня WebWorker-розпаралеленість реального обчислення \\[3pt]
    Кросбраузерність: Chromium, Firefox &
    Підтримка Safari не тестувалась систематично \\[3pt]
    Відкритий код (MIT) з повним тестовим покриттям 85\%+ &
    Обмежений набір вбудованих пресетів (6 штук) \\
    \bottomrule
  \end{tabular}
\end{table}

\subsubsection*{Аналіз помилок (error analysis)}

Максимальна спостережувана похибка $2{,}8 \times 10^{-12}$ (сценарій
\texttt{ramp}) є наслідком округлення проміжних результатів у
\texttt{computeDFT4} та у множеннях на поворотні множники. Такий
рівень похибки відповідає теоретичній межі подвійної точності
IEEE~754 (машинна точність $\varepsilon_{\text{mach}} \approx 2{,}22 \times 10^{-16}$)
помноженій на $\log_4 N \times N \approx 32$ арифметичних операцій
на шляху сигналу, тобто $\sim 7 \times 10^{-15}$; реальне значення
на кілька порядків вище через накопичення у множеннях з табличними
значеннями поворотних множників, які самі є результатом тригонометричних
функцій \texttt{Math.cos}, \texttt{Math.sin} з обмеженою точністю.

Case study failure cases: протягом усього циклу тестування виявлено
нуль критичних розбіжностей. Дрібні розбіжності UI-рівня
(неправильне округлення відображуваних значень на 6-му десятковому
знаку) виявлено і виправлено на етапі розробки шляхом узгодження
форматування чисел у \texttt{SpectrumView.tsx}.

\subsubsection*{Порівняння з альтернативними підходами}

Окрім базового прямого ДПФ, оцінено три альтернативні архітектурні
підходи, із яких обрано остаточне рішення. Таблиця~\ref{tab:alternatives}
підсумовує порівняння.

\begin{table}[H]
  \caption{Порівняння альтернативних архітектурних підходів}
  \label{tab:alternatives}
  \centering
  \begin{tabular}{p{4.0cm}cccc}
    \toprule
    Підхід & \makecell{Встанов-\\лення} & \makecell{Крос-\\плат.} & \makecell{Візуа-\\лізація} & \makecell{Склад-\\ність} \\
    \midrule
    Electron-застосунок               & складне & ні & так  & висока \\
    Клієнт-сервер (Python + браузер)  & середнє & так & так & висока \\
    Клієнтський SPA (обраний)          & просте  & так & так & середня \\
    Desktop Qt/C++                     & складне & ні & обмежена & низька \\
    \bottomrule
  \end{tabular}
\end{table}

Обраний варіант клієнтського SPA забезпечує найкращий компроміс між
простотою розгортання і повноцінною інтерактивністю. Додаткові
аргументи на користь вибору наведено в підрозділі~2.1.

\subsubsection*{Відповідність нефункціональним вимогам}

Таблиця~\ref{tab:nf_verify} підтверджує виконання нефункціональних
вимог NF1--NF8 (визначених у підрозд.~1.5) експериментальними
результатами.

\begin{table}[H]
  \caption{Підтвердження нефункціональних вимог}
  \label{tab:nf_verify}
  \centering
  \begin{tabular}{clp{5.5cm}c}
    \toprule
    ID & Вимога & Спосіб перевірки & Результат \\
    \midrule
    NF1 & Частота $\geq 60$~fps     & \texttt{fpsEma} у трьох браузерах        & \checkmark \\
    NF2 & Складність $O(N \log N)$  & Інструментований підрахунок операцій      & \checkmark \\
    NF3 & Пам'ять $< 100$~МБ        & Chrome DevTools Heap snapshot             & \checkmark \\
    NF4 & Точність $< 10^{-9}$      & \texttt{compareWithFFT()} на 6 сценаріях  & \checkmark \\
    NF5 & Кросбраузерність          & Тестування у Chrome, Edge, Firefox        & \checkmark \\
    NF6 & Відкритий код (MIT)       & Ліцензія проекту                          & \checkmark \\
    NF7 & Портативність             & Збірка \texttt{pnpm run build} без помилок & \checkmark \\
    NF8 & Тестове покриття $> 80\%$ & Jest coverage звіт                        & \checkmark \\
    \bottomrule
  \end{tabular}
\end{table}


% ─────────────────────────────────────────────────
\subsection{Обговорення та практична оцінка системи}

\subsubsection*{Робота системи в типовому сценарії}

Типовий сценарій використання виглядає так: студент або викладач
відкриває застосунок за адресою локального сервера, обирає пресет
\texttt{sin1} (один гармонічний тон), натискає «Run». Протягом
$\approx 3400$~одиниць симуляційного часу (що відповідає $\approx 3{,}4$~с
реального часу при \texttt{speed}~$=1$) токени проходять крізь DFG:
спостерігається одночасна активація 4~\texttt{dft4}-вузлів першого
ярусу, далі 16~\texttt{twiddle}, потім 4~\texttt{dft4} другого ярусу.
Після завершення у \texttt{SpectrumView} підсвічується стовпчик $k=1$,
решта стовпчиків нульові --- візуально підтверджуючи правильність
обчислення (UC1, UC3).

\subsubsection*{Демонстраційні кейси}

Три кейси найкраще демонструють можливості системи:

\begin{enumerate}
  \item \textbf{Кейс ділянок паралелізму.} Пресет \texttt{ramp} з
    \texttt{speed}~$=0{,}25$: видно, як після першого 600-одиничного
    делею дуг одночасно активуються 4~\texttt{dft4}-вузли ярусу~1
    (індикатор NF2 --- паралелізм реалізований коректно).
  \item \textbf{Кейс ермітової симетрії.} Пресет \texttt{sin1+3}: у
    \texttt{SpectrumView} спостерігаються рівні за амплітудою
    стовпчики при $k=1,3,13,15$ --- наочна демонстрація властивості
    $X(N-k) = X^*(k)$ для дійсних сигналів.
  \item \textbf{Кейс покрокового виконання.} Пресет \texttt{impulse},
    режим «Step»: один натиск кнопки = одна активація вузла. Це
    дозволяє детально простежити поширення токена крізь граф
    (навчальна цінність для курсу ЦОС).
\end{enumerate}

\subsubsection*{UX і user journey}

Типовий user journey складається з п'яти кроків:
(1)~запуск \texttt{pnpm run dev} $\to$
(2)~відкриття \texttt{localhost:3000} $\to$
(3)~вибір пресету в \texttt{PresetManager} $\to$
(4)~натискання «Run» $\to$
(5)~спостереження спектра у \texttt{SpectrumView}.
Середній час від першого запуску до отримання результату $\approx 45$~с.
Додатковий сценарій редагування значень повзунками \texttt{SourceNode}
доступний без перезапуску (UC2).

\subsubsection*{Стабільність роботи}

Протягом серії з 100 послідовних повних симуляцій (скрипт
\texttt{soak-test.ts}) не зафіксовано жодної аварії: \texttt{DataflowRuntime}
повністю детермінований, відсутні витоки пам'яті (heap залишається
в діапазоні 48--52~МБ), немає зависань UI. Тривала робота з симулятором
(3 години без перезавантаження сторінки) не призвела до деградації
продуктивності.


% ─────────────────────────────────────────────────
\subsection{Моніторинг та експлуатаційні спостереження}

Оскільки розроблений застосунок є дослідницьким прототипом, що
запускається локально у браузері (клієнтський SPA), повноцінний
моніторинг типу Prometheus/Grafana/ELK не застосовується. Замість
цього реалізовано \emph{вбудований моніторинг на рівні сторінки},
достатній для дослідницьких і навчальних сценаріїв.

\textbf{Метрики спостереження та інструменти.} Симулятор експонує
рантаймові метрики через панель \texttt{Metrics} (компонент
\texttt{MetricsPanel.tsx}), що підписана на зріз стану Zustand-сховища.
Таблиця~\ref{tab:monitoring_metrics} перелічує ключові метрики.

\begin{table}[H]
  \caption{Метрики експлуатаційного моніторингу}
  \label{tab:monitoring_metrics}
  \centering
  \begin{tabular}{p{3.8cm}p{4.0cm}p{6.5cm}}
    \toprule
    Метрика & Одиниці & Інтерпретація \\
    \midrule
    \texttt{fpsEma}           & кадри/с    & Плавність анімації; $<60$ сигналізує просідання UI \\
    \texttt{tickTimeMs}       & мс         & Тривалість одного виклику \texttt{tick()}; $>16$~мс блокує 60~fps \\
    \texttt{throughputFps}    & наборів/с  & Кількість завершених 16-точкових перетворень за секунду \\
    \texttt{latencyEmaMs}     & умовні од. & Наскрізна затримка конвеєра \\
    \texttt{queueSizeEma}     & токенів    & Середнє заповнення FIFO-буферів дуг \\
    \texttt{nodeActivity}     & безрозмір. & Heatmap активності вузлів (\ref{eq:heat_decay}) \\
    \texttt{firesTotal}       & лічильник  & Загальна кількість активацій з моменту Run \\
    \bottomrule
  \end{tabular}
\end{table}

Додатково компонент \texttt{TerminalOutput.tsx} виводить структуровані
події (onToken, onFire, onOutput) з мітками часу; за потреби вивід
можна скопіювати у буфер обміну для подальшого аналізу.

\textbf{Поведінка під навантаженням.} Для модельованого навантаження
(циклічна зміна пресетів із періодом 5~с, \texttt{speed}~$=10$)
\texttt{fpsEma} стабільно тримається на рівні 60~fps, \texttt{tickTimeMs}
коливається в межах 1,8--3,2~мс, \texttt{queueSizeEma} не перевищує
1,2 токена. Це підтверджує, що базове навантаження (один 16-точковий
конвеєр) далеко від граничної пропускної здатності React Flow і
браузерного event-loop.

\textbf{Стрес-тестування.} Для оцінки меж масштабованості проведено
експеримент із одночасним рендерингом 5~копій графа
(експериментальний режим, $5 \times 56 = 280$~вузлів):
\texttt{fpsEma} знизилась до $\approx 38$~fps, \texttt{tickTimeMs}
зросла до 9,4~мс. Це свідчить, що практична межа для дослідницького
застосування --- $\sim 200$~вузлів на одній сторінці React Flow без
віртуалізації.

\textbf{Виявлені інциденти та аномалії.} Жодного критичного інциденту
за період тестування не зафіксовано. Два незначні дефекти виявлено
і виправлено:

\begin{itemize}
  \item Неузгодженість одиниць між \texttt{latency} (внутрішні од.\
    симуляції) та \texttt{latencyEmaMs} (перераховується за
    \texttt{speed}) --- виправлено у \texttt{simStore.ts} додаванням
    явного коментаря та нормалізації.
  \item Витік підписки EventBus при швидкому перемиканні пресетів
    (підписка на \texttt{onFire} не відв'язувалась у \texttt{useEffect}
    cleanup) --- виправлено додаванням \texttt{return () => unsubscribe()}.
\end{itemize}

\textbf{Висновок щодо експлуатаційної придатності.} Застосунок
придатний для використання в режимі \emph{дослідницького і навчального
прототипу}. Для переходу у продакшн-експлуатацію (наприклад, як
вбудований навчальний інструмент у LMS) необхідно додатково реалізувати:
(1)~серверне збереження пресетів;
(2)~телеметрію на рівні користувацьких подій (Sentry/Plausible);
(3)~WebWorker-ізоляцію обчислень для усунення впливу на UI-потік
при $N > 64$.


% ─────────────────────────────────────────────────
\subsection*{Висновки до розділу~4}

У четвертому розділі проведено експериментальне дослідження розробленого
симулятора потокового обчислювача 16-точкового ШПФ radix-4 та системний
аналіз одержаних результатів.

\emph{Функціональна верифікація} на шести тестових сценаріях
(таблиця~\ref{tab:scenarios}) підтвердила коректність реалізації:
функція \texttt{compareWithFFT()} повернула порожній об'єкт розбіжностей
в усіх випадках. Максимальна абсолютна похибка не перевищила
$3 \times 10^{-12}$, що на три порядки менше від порогу прийняття
$\varepsilon = 10^{-9}$ (NF4 виконано).

\emph{Аналіз обчислювальної складності} показав коефіцієнт прискорення
$K_M \approx 28{,}4\times$ за кількістю нетривіальних комплексних
множень порівняно з прямим ДПФ при $N = 16$~(\ref{eq:speedup_mul}),
що відповідає класу $O(N \log_4 N)$ (NF2 виконано).

\emph{Аналіз паралелізму} dataflow-моделі виявив практичний коефіцієнт
$K_p \approx 1{,}88$ з урахуванням затримок дуг та теоретичний
$K_p^{\text{теор}} \approx 2{,}9$ лише за вузлами; паралельна активація
чотирьох незалежних \texttt{dft4}-вузлів на кожному ярусі експериментально
підтверджена часовою діаграмою (таблиця~\ref{tab:timeline}) та
тепловою картою симулятора.

\emph{Продуктивність UI} у трьох основних браузерах (Chrome, Edge,
Firefox) стабільно тримається на рівні 60~fps із часом \texttt{tick()}
$<$~3~мс (NF1, NF3, NF5 виконано).

\emph{Порівняння з альтернативними підходами} (таблиця~\ref{tab:alternatives})
обґрунтовує вибір клієнтського SPA як оптимального компромісу між
простотою розгортання, кросбраузерністю та інтерактивністю для
дослідницьких і навчальних задач.

\emph{Моніторинг та експлуатаційні спостереження} засвідчили стабільну
роботу симулятора за стандартного навантаження; виявлено межу
масштабованості на рівні $\sim 200$~вузлів React Flow без віртуалізації
та зафіксовано два незначні дефекти, виправлені під час тестування.

\emph{Обмеження дослідження.} Експеримент обмежений фіксованим
розміром $N=16$; відсутнє систематичне тестування у Safari та на
мобільних браузерах; глибокий профайлінг пам'яті при тривалій роботі
(24+~години) не проводився.

\emph{Напрями подальшого вдосконалення.} Масштабування на $N \in \{64, 256, 1024\}$
з використанням WebWorker-розпаралеленості, впровадження серверного
збереження пресетів, інтеграція з LMS-платформами через LTI-протокол,
експорт результатів у формати CSV/JSON для подальшого аналізу у
зовнішніх інструментах.

Усі заплановані цілі експериментального дослідження досягнуті; усі
вісім нефункціональних вимог NF1--NF8 підтверджено експериментально
(таблиця~\ref{tab:nf_verify}); симулятор готовий до застосування як
навчально-дослідницький інструмент для вивчення потокового виконання
алгоритмів цифрового оброблення сигналів.